\chapter{The Mismatch String Kernel}\label{ch:kernel}

\section{Kernel Methods: A Brief Introduction}\label{sec:kernel-intro}

In the field of machine learning, assessing the similarity between data points is a fundamental requirement. However, when dealing with complex, non-vectorial data such as discrete strings or genomic sequences, defining this similarity is not straightforward. Kernel methods provide an elegant mathematical framework to solve this by bridging the gap between abstract objects and structured linear spaces.

\subsection{The Kernel Trick and Feature Spaces}\label{sec:kernel-trick}

Many standard learning algorithms, such as Support Vector Machines (SVMs) that will be used in this work, rely on dot products to evaluate similarity between patterns. However, input patterns do not necessarily exist in a native dot product space initially; they could be any kind of discrete object, including biological sequences.

To apply linear analytical geometry to these objects, we must first map them from their original input domain $\mathcal{X}$ into a dot product space $\mathcal{H}$, commonly referred to as the feature space. This transformation is performed via a mapping function $\Phi$:

\begin{equation}\label{eq:feature-map}
x \mapsto \mathbf{x} := \Phi(x)
\end{equation}

\begin{definizione}[Kernel Function]\label{def:kernel}
Embedding data into the feature space $\mathcal{H}$ via $\Phi$ allows us to define a similarity measure directly from the dot product in that space. A function $k$ that computes this similarity is called a kernel, defined as:
\begin{equation}\label{eq:kernel-def}
k(x, x') := \langle \mathbf{x}, \mathbf{x'} \rangle = \langle \Phi(x), \Phi(x') \rangle
\end{equation}
\end{definizione}

This substitution is known as the ``kernel trick''. The power of the kernel trick lies in its ability to compute the inner product $\langle \Phi(x), \Phi(x') \rangle$ directly in the input space $\mathcal{X}$ without ever explicitly constructing or evaluating the potentially high-dimensional (or even infinite-dimensional) feature vectors $\Phi(x)$ and $\Phi(x')$. Given an algorithm formulated entirely in terms of a positive definite kernel $k$, one can effortlessly replace $k$ with another positive definite kernel, effectively changing the feature space mapping without altering the underlying dot product-based algorithm.

\subsection{Theoretical Foundations (RKHS and Mercer's Theorem)}\label{sec:rkhs}

For a function $k(x, x')$ to be considered a valid kernel that computes a dot product in some feature space $\mathcal{H}$, it must satisfy specific mathematical properties. Primarily, the kernel must be symmetric and give rise to a positive definite Gram matrix for any set of input patterns.

Any positive definite kernel can be associated with a Reproducing Kernel Hilbert Space (RKHS). In this functional space, the kernel acts as the representer of evaluation. The dot ($\cdot$) in the notation $k(x, \cdot)$ acts as a placeholder for a free variable: while the first argument is fixed to a specific data point $x$, the second argument is left open, transforming the two-variable kernel into a function of just one variable. For any function $f \in \mathcal{H}$, the reproducing property holds:

\begin{equation}\label{eq:reproducing-property}
\langle f, k(x, \cdot) \rangle = f(x)
\end{equation}

This equation states that taking the inner product of any function $f$ in the Hilbert space $\mathcal{H}$ with the partially-fixed kernel function $k(x, \cdot)$ results simply in the function $f$ evaluated at the point $x$. 
The kernel effectively ``reproduces'' the evaluation of the function. 

Consequently, by substituting the generic function $f$ with another kernel function $k(x', \cdot)$, taking the dot product of the kernel with itself yields the kernel function evaluated at those original points:

\begin{equation}\label{eq:kernel-dot}
\langle k(x, \cdot), k(x', \cdot) \rangle = k(x, x')
\end{equation}

This is the most powerful consequence of the reproducing property and the foundation of the kernel trick. 
It guarantees that the inner product of two feature mappings in the high-dimensional Hilbert space is exactly equal to the original kernel function evaluated on the original points $x$ and $x'$. 

The theoretical justification for why such feature spaces exist is provided by Mercer's Theorem. 
\begin{teorema}[Mercer's Theorem]\label{thm:mercer}
A symmetric, real-valued function $k \in L_{\infty}(\mathcal{X}^2)$ corresponds to a dot product in a feature space if the integral operator associated with it is positive definite. Specifically, for all functions $f \in L_2(\mathcal{X})$, the kernel must satisfy the condition:
\begin{equation}\label{eq:mercer-condition}
\int_{\mathcal{X}^2} k(x, x') f(x) f(x') \, d\mu(x) d\mu(x') \geq 0
\end{equation}
\end{teorema}

Because Mercer kernels are positive definite, they are also reproducing kernels, allowing them to be seamlessly represented as dot products within Hilbert spaces. This mathematical foundation guarantees that as long as a function is a valid positive definite kernel, we can safely compute similarity measures for complex structures like genomic sequences, knowing that a corresponding feature space $\mathcal{H}$ implicitly exists to support the operations.

\section{String Kernels for Genomic Sequences}\label{sec:string-kernels}

Standard continuous kernels, such as the Radial Basis Function (RBF), are typically designed for continuous vectorial data and often fail to capture the underlying structure of discrete biological sequences. Genomic sequences are intrinsically discrete, vary in length, and lack a natural geometric alignment in a standard Euclidean space. To address this, string kernels map these sequences into a high-dimensional feature space based on their sub-string occurrences. In this work, the sequences are composed of an extended epigenetic alphabet:
\begin{equation}\label{eq:alphabet}
\Sigma = \{A, C, G, T, M, H\}
\end{equation}
which incorporates both standard nucleotides and specific epigenetic modifications.

\subsection{The Spectrum Kernel}\label{sec:spectrum}

The simplest approach to defining a sequence similarity measure is through the exact matching of short subsequences, known as $k$-mers. This methodology forms the basis of the spectrum kernel, which constructs a feature space indexed by all possible sequences of length $k$. 

\begin{definizione}[Spectrum Feature Map]\label{def:spectrum-feature-map}
The spectrum feature map for a sequence $x$ at k-mer length $k$ is defined as the vector of k-mer counts:
\begin{equation}\label{eq:spectrum-feature-map}
\Phi_k(x) = (\phi_a(x))_{a \in \mathcal{A}^k}
\end{equation}
\end{definizione}

In this representation, each coordinate corresponds to the number of times a specific $k$-mer appears in the input sequence. Because any given sequence of length $n$ contains at most $n - k + 1$ distinct $k$-mers, the resulting feature vector is highly sparse, even though the theoretical dimensionality of the space ($|\Sigma|^k$) can be exceptionally large.

\paragraph{Example (Spectrum Feature Map)}
Consider a simplified alphabet $\Sigma = \{\text{A}, \text{C}, \text{G}, \text{T}\}$ and $k=3$. Let $x = \text{ACGAG}$ be a short sequence. 
The sequence contains exactly three overlapping $3$-mers: \texttt{ACG}, \texttt{CGA}, and \texttt{GAG}.
Therefore, the spectrum feature map $\Phi_3(x)$ will be a high-dimensional vector where the coordinates corresponding to \texttt{ACG}, \texttt{CGA}, and \texttt{GAG} each have a count of $1$. All other $4^3 - 3 = 61$ dimensions in the feature space evaluate to $0$. 
If another sequence $y$ shares none of these exact $3$-mers, their inner product $\langle \Phi_3(x), \Phi_3(y) \rangle$ will evaluate to strictly zero.

\begin{definizione}[Spectrum Kernel]\label{def:spectrum-kernel}
The spectrum kernel is the inner product of the spectrum feature representations:
\begin{equation}\label{eq:spectrum-kernel}
K_k(x, y) = \langle \Phi_k(x), \Phi_k(y) \rangle
\end{equation}
\end{definizione}

While the spectrum kernel is computationally efficient, it exhibits limitations in biological applications. Since it relies exclusively on exact matches, a single nucleotide variation or sequencing error completely alters the $k$-mer, preventing the kernel from recognizing the structural similarity between two closely related motifs.

\subsection{The Mismatch Kernel}\label{sec:mismatch-kernel}

To overcome the limitations of exact matching, the mismatch kernel introduces a degree of flexibility that accounts for natural variations and sequencing noise without requiring complex generative models, such as Hidden Markov Models. Instead of matching $k$-mers strictly, the mismatch kernel considers a neighborhood of similar sequences.

\begin{definizione}[Mismatch Neighbourhood]\label{def:mismatch-neighbourhood}
The $(k,m)$-mismatch neighbourhood of a k-mer $\alpha$ is the set of all k-mers within Hamming distance $m$:
\begin{equation}\label{eq:mismatch-neighbourhood}
N_{(k,m)}(\alpha) = \{ \beta \in \Sigma^k : d_H(\alpha, \beta) \leq m \}
\end{equation}
\end{definizione}

\begin{definizione}[Mismatch Feature Map]\label{def:mismatch-feature-map}
The mismatch feature map for a single k-mer $\alpha$ assigns a weight of 1 to all k-mers within the mismatch neighbourhood:
\begin{equation}\label{eq:mismatch-feature-kmer}
\Phi_{(k,m)}(\alpha) = (\phi_{\beta}(\alpha))_{\beta \in \Sigma^k}
\end{equation}
The feature map for a full sequence $x$ is obtained by summing over all k-mers in $x$:
\begin{equation}\label{eq:mismatch-feature-seq}
\Phi_{(k,m)}(x) = \sum_{k\text{-mers } \alpha \text{ in } x} \Phi_{(k,m)}(\alpha)
\end{equation}
\end{definizione}

By expanding the feature representation additively over the mismatch neighborhood, a single $k$-mer observation contributes not only to its exact corresponding coordinate but also to the coordinates of all its $m$-mismatch neighbors. Consequently, sequences that share homologous but not identical motifs will still produce a significant inner product.

\paragraph{Example (Mismatch Neighbourhood)}
Following the previous example with $\Sigma = \{\text{A}, \text{C}, \text{G}, \text{T}\}$ and $k=3$, suppose we allow a mismatch tolerance of $m=1$. 
Let us consider the single $3$-mer $\alpha = \text{\texttt{ACG}}$. Its $(3,1)$-mismatch neighbourhood $N_{(3,1)}(\text{\texttt{ACG}})$ consists of the exact sequence itself and all sequences that differ by exactly one character. Specifically, it includes:
\begin{itemize}
    \item \textbf{Distance $0$ (Exact match):} \texttt{ACG}
    \item \textbf{Distance $1$ (Position 1 changed):} \texttt{CCG}, \texttt{GCG}, \texttt{TCG}
    \item \textbf{Distance $1$ (Position 2 changed):} \texttt{AAG}, \texttt{AGG}, \texttt{ATG}
    \item \textbf{Distance $1$ (Position 3 changed):} \texttt{ACA}, \texttt{ACC}, \texttt{ACT}
\end{itemize}
In the mismatch feature map $\Phi_{(3,1)}(x)$, an observation of \texttt{ACG} in the sequence will result in a $+1$ increment for all 10 of these coordinates. 
If another sequence contains the near-match \texttt{ACT}, the two sequences will share multiple active coordinates in their feature vectors, resulting in a positive similarity score despite lacking an exact $k$-mer match.

\subsection{Weighted Mismatch Neighbourhood}\label{sec:weighted-mismatch}

The standard formulation of the mismatch kernel employs a strict binary cutoff: neighbors within distance $m$ receive a weight of $1$, and those beyond receive $0$. While effective, this step-function approach treats a perfect match identically to a match with $m$ errors. To refine this representation, we introduce a continuous decay penalty based on the exact distance.

\begin{definizione}[Weighted Mismatch]\label{def:weighted-mismatch}
Rather than a binary cutoff, mismatch neighbours are penalized by their exact Hamming distance through a decay factor:
\begin{equation}\label{eq:weighted-mismatch}
w(\beta, \alpha) = \lambda^{d_H(\alpha, \beta)} \quad \text{where} \quad \lambda \in [0, 1]
\end{equation}
\end{definizione}

This modification ensures a smoother feature space. By weighting the contributions, we mathematically enforce that exact biological matches contribute the maximum similarity score, while near-matches provide a geometrically decreasing degree of support.

\section{Multiple Kernel Learning (MKL)}\label{sec:mkl}

Genomic sequences harbor information across multiple scales. Short $k$-mers might effectively capture localized methylation states or localized transcription factor binding affinities, whereas longer $k$-mers are better suited for identifying wider structural motifs. Relying on a single static value of $k$ limits the classifier's capacity. Multiple Kernel Learning (MKL) addresses this by linearly combining kernels computed at various lengths.

\begin{definizione}[MKL Combination]\label{def:mkl}
The multiple kernel learning combination aggregates kernels across a range of k-mer lengths with weights $w_k$:
\begin{equation}\label{eq:mkl}
K_{\text{MKL}}(x, x') = \sum_{k=k_{\min}}^{k_{\max}} w_k \, K_k(x, x')
\end{equation}
\end{definizione}

To optimize this aggregation, the weighting parameters $w_k$ can be structured as a linear ramp. This strategy artificially suppresses the contribution of excessively short $k$-mers, which are highly abundant and often act as background noise, while emphasizing longer, more uniquely discriminative biological patterns.

The starting point of this ramp ($k_{\min}$) must also be logically constrained by the mismatch tolerance $m$. For example, permitting $m=1$ on a $2$-mer effectively ignores half of the subsequence's informational content. This disproportionate error allowance would flood the feature space with false-positive alignments, negating the utility of tracking short motifs entirely.

\section{Gram Matrix Computation and Normalization}\label{sec:gram}

The execution of kernel methods within a classification or novelty detection framework requires the explicit evaluation of the Gram matrix, which encapsulates all pairwise similarities. 

\subsection{Symmetric Gram Matrix (Training)}\label{sec:gram-symmetric}

During the training phase, the core operation involves mapping the relationships between all available training sequences. 

The Gram matrix entry for training pairs is computed as:
\begin{equation}\label{eq:gram-entry}
K_{ij} = \langle \Phi(x_i), \Phi(x_j) \rangle
\end{equation}

Because the kernel function evaluates a standard inner product, the resulting matrix is inherently symmetric ($K_{ij} = K_{ji}$). This symmetry halves the computational load, as only the upper triangle and the diagonal need to be computed. 

Furthermore, raw inner products scale with the length of the sequences being compared. Without correction, longer sequences will artificially register higher similarity scores simply by virtue of generating more $k$-mers. To ensure that similarity relies purely on compositional overlap, we apply cosine normalization.

The kernel is normalized to correct for varying sequence lengths:
\begin{equation}\label{eq:gram-normalize}
\hat{K}_{ij} = \frac{K_{ij}}{\sqrt{K_{ii} K_{jj}}}
\end{equation}

\subsection{Asymmetric Gram Matrix (Inference)}\label{sec:gram-asymmetric}

Inference requires determining the relationship between an unseen test distribution and the established training support vectors. Unlike the square, symmetric training matrix, the evaluation matrix is generally asymmetric.

The test Gram matrix computes similarities between test and training sequences:
\begin{equation}\label{eq:gram-test}
K_{\text{test}}[i, j] = \langle \Phi(x_i^{\text{test}}), \Phi(x_j^{\text{train}}) \rangle
\end{equation}

A fundamental challenge at this stage is aligning the feature space vocabulary between the independent sequence distributions. Furthermore, applying the proper cosine normalization at inference requires access to the individual geometric lengths of the test sequences in feature space.

Proper normalization at test time utilizes test self-norms alongside training self-norms:
\begin{equation}\label{eq:gram-test-normalize}
\hat{K}_{\text{test}}[i, j] = \frac{K_{\text{test}}[i, j]}{\sqrt{K^{\text{test}}_{ii} K^{\text{train}}_{jj}}}
\end{equation}

By independently calculating the self-norm diagonal for the test set ($K^{\text{test}}_{ii}$), we guarantee that the classifier's spatial decision boundaries are respected, preventing test sequences of differing lengths from systematically skewing the anomaly detection margins.

\section{Kernel Approximation via the Nystr\"{o}m Method}\label{sec:nystrom-theory}

\subsection{The Computational Bottleneck}\label{sec:nystrom-bottleneck}

While the exact Gram matrix provides a flawless representation of the geometric similarities in the feature space, its explicit construction introduces a severe computational bottleneck. For a dataset of $N$ sequences, calculating the exact symmetric Gram matrix requires computing and storing $N(N+1)/2$ unique dot products. Both the time complexity and the memory footprint scale quadratically, $O(N^2)$. 

For the massive datasets typical in genomic studies---where liquid biopsy samples can yield millions of reads---computing the exact matrix becomes intractable. For example, storing an exact Gram matrix for $100\,000$ sequences in 64-bit floating-point precision demands approximately $80$\,GB of RAM. Consequently, to utilize robust kernel methods on such large cohorts without requiring highly specialized hardware, we must employ a low-rank approximation strategy.

\subsection{Mathematical Formulation}\label{sec:nystrom-math}

The Nystr\"{o}m method is a technique used to generate a low-rank approximation of a large positive semi-definite Gram matrix. The core intuition relies on selecting a small subset of training samples to act as a representative basis, rather than computing all pairwise similarities.

Let $X$ be the full dataset of $N$ samples. We select $m$ samples uniformly at random to serve as ``landmark'' points, where $m \ll N$. These landmarks form the subset $X_m$. We can then define two sub-matrices based on kernel evaluations:
\begin{itemize}
    \item $W$: The $m \times m$ landmark kernel matrix, where $W_{ij} = K(x_i, x_j)$ for $x_i, x_j \in X_m$.
    \item $C$: The $N \times m$ cross-kernel matrix, where $C_{ij} = K(x_i, x_j)$ for $x_i \in X$ and $x_j \in X_m$.
\end{itemize}

The standard Nystr\"{o}m approximation $\tilde{K}$ for the full Gram matrix $K$ is given by:
\begin{equation}\label{eq:nystrom-approx}
\tilde{K} = C W^{-1} C^T
\end{equation}

By utilizing this approximation, the computational complexity drops from $O(N^2)$ to $O(N \cdot m + m^3)$, effectively linearizing the cost with respect to the total dataset size $N$, since $m$ is relatively small.

To apply this approximation within our Support Vector Machine framework, it is advantageous to extract an explicit, low-dimensional dense feature map $\Phi_{\text{Nys}} \in \mathbb{R}^{N \times m}$ such that $\tilde{K} = \Phi_{\text{Nys}} \Phi_{\text{Nys}}^T$.
Because $W$ is positive semi-definite, we can compute its eigendecomposition:
\begin{equation}\label{eq:w-eigendecomp}
W = V \Lambda V^T
\end{equation}
where $V$ contains the eigenvectors and $\Lambda$ is the diagonal matrix of eigenvalues. We can then compute the inverse square root of $W$:
\begin{equation}\label{eq:w-inv-sqrt}
W^{-1/2} = V \Lambda^{-1/2} V^T
\end{equation}

The explicit feature projection for the entire dataset is defined as:
\begin{equation}\label{eq:nystrom-feature-map}
\Phi_{\text{Nys}} = C W^{-1/2}
\end{equation}

This mapping projects the highly sparse string kernel representation into a dense, $m$-dimensional Euclidean space. Consequently, instead of solving a non-linear kernel SVM using the intractable $N \times N$ Gram matrix, we can train a significantly faster linear Support Vector Machine directly on the dense features $\Phi_{\text{Nys}}$, while theoretically maintaining the non-linear decision boundary defined by the mismatch string kernel.

\subsection{Measuring Approximation Fidelity}\label{sec:nystrom-metrics-theory}

Because the Nystr\"{o}m method is a low-rank approximation, it inherently introduces an error matrix $E = K - \tilde{K}$. To quantify the quality of the approximation and validate the chosen number of landmarks $m$, we define three intrinsic mathematical metrics:

\begin{enumerate}
    \item \textbf{Relative Frobenius Error:} This measures the global element-wise deviation across the entire matrix. It is defined using the Frobenius norm:
    \begin{equation}\label{eq:frob-error}
    \epsilon_F = \frac{\|K - \tilde{K}\|_F}{\|K\|_F} = \frac{\sqrt{\sum_{i,j} (K_{ij} - \tilde{K}_{ij})^2}}{\|K\|_F}
    \end{equation}

    \item \textbf{Relative Spectral Error:} This evaluates the maximum distortion introduced to the geometric structure of the feature space, computed via the spectral norm ($L_2$ operator norm), which corresponds to the largest singular value of the error matrix:
    \begin{equation}\label{eq:spectral-error}
    \epsilon_2 = \frac{\|K - \tilde{K}\|_2}{\|K\|_2}
    \end{equation}

    \item \textbf{Diagonal Error:} The exact mismatch kernel relies on cosine normalization, ensuring that the self-similarity for every sequence is exactly $1$ ($K_{ii} = 1$). The diagonal error evaluates how well these self-similarities (or sequence lengths in the feature space) are preserved by the approximation. It is captured by the maximum and mean deviations:
    \begin{equation}\label{eq:diag-error}
    \epsilon_{\text{diag-max}} = \max_i |K_{ii} - \tilde{K}_{ii}|, \quad \epsilon_{\text{diag-mean}} = \frac{1}{N} \sum_{i=1}^N |K_{ii} - \tilde{K}_{ii}|
    \end{equation}
\end{enumerate}

Minimizing these intrinsic errors guarantees that the explicit dense projection $\Phi_{\text{Nys}}$ remains geometrically faithful to the exact mismatch string kernel space.